\documentclass[../../main.tex]{subfiles}% 注意这里的文件路径不能用 ./main.tex ,否则用latexmk编译子文件会报错
\graphicspath{{\subfix{./image/}}{\subfix{../../image/}}} % 指定图片引用路径,后续可以直接使用图片文件名
% 如果图片存放在上级目录,则使用 ../../image/ 作为备用路径

% 例如:
% \begin{figure}[H]
% \centering
% \includegraphics[width=0.3\textwidth]{图.png}
% \caption{}
% \label{figure:图}
% \end{figure}
% 注意:上述\label{}一定要放在\caption{}之后,否则引用图片序号会只会显示??.

\begin{document}

\section{存在性}\label{section:ch08-存在性}

借助上节结果，可证明如下存在性定理。

\begin{theorem}[Young 测度值解的整体存在性]\label{theorem:ch08-存在性}
设或 $N=1$ 且\eqref{align:ch08-8.1a}--\eqref{align:ch08-8.1d} 成立，或 $N\geqslant2$ 且\eqref{align:ch08-8.2a}--\eqref{align:ch08-8.2e} 成立。则问题 (P) 存在整体 Young 测度值解 $(u,\nu)$，它作为正则化问题族 $(P\epsilon_k)$ 解的某子列 $\{u_{\epsilon_k}\}$ 的极限而得到。进而：
\begin{enumerate}[(i)]
\item\label{item:ch08-exist-i} 若 $N=1$，则 $u\in L^\infty(0,\infty;W^{1,p}_0(\Omega))\cap C([0,\infty);L^2(\Omega))\cap H^1(Q_T)$（对所有 $p\in[1,\infty)$），且 $u_x\in L^\infty(Q_\infty)$；
\item\label{item:ch08-exist-ii} 若 $N\geqslant2$，则 $u\in L^\infty(0,\infty;W^{1,p}_0(\Omega))\cap C([0,\infty);L^p(\Omega))\cap W^{1,p}(Q_T)$（$p\in(1,2]$）；
\item\label{item:ch08-exist-iii} $\nu\in\mathcal{Y}^+(Q_\infty;\mathbb{R}^N)$。
\end{enumerate}
\end{theorem}

\begin{proof}
先设 $N\geqslant2$。考察\cref{proposition:ch08-收敛N2} 给出的序列 $\{u_{\epsilon_k}\}$ 与函数 $u$。问题 $(P\epsilon_k)$ 的弱形式为
\begin{align}
&\iint_{Q_t}[u_{\epsilon_k}\partial_s\zeta-\phi(\nabla u_{\epsilon_k})\cdot\nabla\zeta-\epsilon_k\nabla u_{\epsilon_k,s}\cdot\nabla\zeta](x,s)\,\mathrm{d}x\,\mathrm{d}s\nonumber\\
&\qquad=\int_\Omega u_{\epsilon_k}(x,t)\zeta(x,t)\,\mathrm{d}x-\int_\Omega u_{0\epsilon_k}(x)\zeta(x,0)\,\mathrm{d}x \label{align:ch08-8.48}
\end{align}
对所有 $t>0$ 与 $\zeta\in C^1([0,\infty);C^1_c(\Omega))$。由\eqref{align:ch08-8.36a}，当 $k\to\infty$ 时对所有 $t>0$ 有
\begin{align}
\lim_{k\to\infty}\iint_{Q_t}u_{\epsilon_k}(x,s)\partial_s\zeta(x,s)\,\mathrm{d}x\,\mathrm{d}s
=\iint_{Q_t}u(x,s)\partial_s\zeta(x,s)\,\mathrm{d}x\,\mathrm{d}s, \label{align:ch08-8.49}\\
\lim_{k\to\infty}\int_\Omega u_{\epsilon_k}(x,t)\zeta(x,t)\,\mathrm{d}x
=\int_\Omega u(x,t)\zeta(x,t)\,\mathrm{d}x, \label{align:ch08-8.50}
\end{align}
而由\eqref{align:ch08-8.36d} 与\eqref{align:ch08-8.28} 有
\begin{align}
\lim_{k\to\infty}\left(\epsilon_k\iint_{Q_t}\nabla\partial_s u_{\epsilon_k}\cdot\nabla\zeta\,\mathrm{d}x\,\mathrm{d}s\right)=0, \label{align:ch08-8.51}
\end{align}
以及
\begin{align}
\lim_{k\to\infty}\int_\Omega u_{0\epsilon_k}(x)\zeta(x,0)\,\mathrm{d}x
=\int_\Omega u_0(x)\zeta(x,0)\,\mathrm{d}x. \label{align:ch08-8.52}
\end{align}
由于由不等式\eqref{align:ch08-8.29} 序列 $\{|\nabla u_{\epsilon_k}|\}$ 在 $L^p(Q_T)$（$p\in(1,2]$）中有界，由\cref{lemma:ch02-有界一致可积} 它一致可积；进而序列 $\{\phi(\nabla u_{\epsilon_k})\}$ 亦然，因由假设\eqref{align:ch08-8.2b} 有 $|\phi(\nabla u_{\epsilon_k})|\leqslant L|\nabla u_{\epsilon_k}|+|\phi(0)|$。于是由\cref{proposition:ch05-弱收敛} 与标准对角论证，存在 $\{\nabla u_{\epsilon_k}\}$ 的子列（为简便不重新标号）与测度 $\nu\in\mathscr{R}_f(Q_\infty\times\mathbb{R}^N)$（对所有 $T\in(0,\infty)$ 有 $\nu\in\mathcal{Y}^+(Q_T;\mathbb{R}^N)$），使
\begin{align}
\phi(\nabla u_{\epsilon_k})\rightharpoonup\phi^*\ \text{在 }L^1(Q_T;\mathbb{R}^N), \label{align:ch08-8.53}
\end{align}
其中 $\phi^*$ 由\eqref{align:ch08-8.6} 定义。在\eqref{align:ch08-8.48} 中令 $k\to\infty$ 并用\eqref{align:ch08-8.49}--\eqref{align:ch08-8.53}，即得等式\eqref{align:ch08-8.5}。类似地，对 $h=\rho_i$（$\rho_i(y)\triangleq y_i$）应用\cref{proposition:ch05-弱收敛}，有
\begin{align*}
(\nabla u_{\epsilon_k})_i\rightharpoonup\int_{\mathbb{R}^N}y_i\,\mathrm{d}\nu_{(x,t)}(y)\ \text{在 }L^1(Q_T)\quad(i=1,\cdots ,N),
\end{align*}
结合\eqref{align:ch08-8.42b} 中的收敛即得等式\eqref{align:ch08-8.4}。

因而偶 $(u,\nu)$ 是问题 (P) 在 $Q_T$ 中的 Young 测度值解。由\cref{proposition:ch08-收敛N2} 得 $u\in L^\infty(0,\infty;W^{1,p}_0(\Omega))\cap C([0,\infty);L^p(\Omega))$ 且 $\partial_t u\in L^2(Q_\infty)$。此外，由于对每个 $T\in(0,\infty)$ 有 $\nu\in\mathcal{Y}^+(Q_T;\mathbb{R}^N)$，由测度的初等性质，对所有 Borel 集 $E\subseteq Q_\infty$ 有
\begin{align*}
\lambda_{N+1}(E)=\lim_{k\to\infty}\lambda_{N+1}(E\cap Q_k)=\lim_{k\to\infty}\nu((E\cap Q_k)\times\mathbb{R}^N)=\nu(E\times\mathbb{R}^N),
\end{align*}
故 $\nu\in\mathcal{Y}^+(Q_\infty;\mathbb{R}^N)$。

由于 $(u,\nu)$ 对所有 $T\in(0,\infty)$ 都是 $Q_T$ 中问题 (P) 的 Young 测度值解，它是整体解；$N\geqslant2$ 情形结论即得。$N=1$ 情形的证明相同，只是用不等式\eqref{align:ch08-8.21} 与\cref{proposition:ch08-收敛N1} 分别代替\eqref{align:ch08-8.29} 与\cref{proposition:ch08-收敛N2}。证毕。
\end{proof}

\begin{remark}\label{remark:ch08-w等于phi星}
设 $N=1$，$\{u_{\epsilon_k}\}$ 为\cref{theorem:ch08-存在性} 证明中所用序列。由不等式\eqref{align:ch08-8.25}，对所有 $T\in(0,\infty)$ 有
\begin{align*}
\|\phi(\partial_x u_{\epsilon_k})-w_{\epsilon_k}\|_{L^2(Q_T)}=\epsilon_k\|\partial_{xt}u_{\epsilon_k}\|_{L^2(Q_T)}\to0,
\end{align*}
故由\eqref{align:ch08-8.45} 有 $\phi(\partial_x u_{\epsilon_k})\rightharpoonup w$ 在 $L^2(Q_T)$。另一方面，由\eqref{align:ch08-8.53} 有 $\phi(\partial_x u_{\epsilon_k})\rightharpoonup\phi^*\triangleq \int_{\mathbb{R}}\phi(y)\,\mathrm{d}\nu(y)$ 在 $L^1(Q_T)$。于是 $w=\phi^*$，特别地 $\phi^*\in L^2(0,\infty;H^1(\Omega))\cap L^\infty(Q_T)$。
\end{remark}

\begin{remark}\label{remark:ch08-phi星可积}
仿\cref{theorem:ch08-存在性} 证明可知，若 $N\geqslant2$ 且假设\eqref{align:ch08-8.2c} 满足，则由\eqref{align:ch08-8.6} 定义的每个函数 $\phi_i^*$ 属于 $L^r(Q_T)$（对所有 $T\in(0,\infty)$ 与 $r=\frac{p}{p-1}$，$p\in(1,2]$；显然 $N=1$ 时由\eqref{align:ch08-8.21} 有 $\phi^*\in L^\infty(Q_\infty)$）。事实上，固定 $T\in(0,\infty)$。由\eqref{align:ch08-8.2c} 存在 $M>0$，使对几乎所有 $(x,t)\in Q_T$，
\begin{align}
|\phi_i^*|(x,t)\leqslant M\int_{\mathbb{R}^N}(1+|y|^{p-1})\,\mathrm{d}\nu_{(x,t)}(y). \label{align:ch08-8.54}
\end{align}
由\eqref{align:ch08-8.54} 与 Jensen 不等式有
\begin{align}
|\phi_i^*|^r\leqslant\bar M\int_{\mathbb{R}^N}(1+|y|^p)\,\mathrm{d}\nu_{(x,t)}(y) \label{align:ch08-8.55}
\end{align}
其中 $\bar M>0$。令 $f_j\in C_c([0,\infty))$（$j\in\mathbb{N}$）满足 $0\leqslant f_j\leqslant1$ 且
\begin{align}
f_j(z)=
\begin{cases}
1 & \text{若 }z\in[0,j],\\
0 & \text{若 }z\in[j+1,\infty).
\end{cases} \label{align:ch08-8.56}
\end{align}
设 $\{u_{\epsilon_k}\}$ 为\cref{theorem:ch08-存在性} 证明中所用序列。由\eqref{align:ch08-8.29}，对每个 $j\in\mathbb{N}$，序列 $\{f_j(|\nabla u_{\epsilon_k}|)|\nabla u_{\epsilon_k}|^p\}$ 在 $Q_T$ 中有界因而一致可积。于是由\cref{proposition:ch05-弱收敛} 存在子列（不重新标号）使
\begin{align}
\iint_{Q_T}\mathrm{d}x\,\mathrm{d}t\int_{\mathbb{R}^N}(1+f_j(|y|)|y|^p)\,\mathrm{d}\nu_{(x,t)}(y)
&=\lim_{k\to\infty}\iint_{Q_T}(1+f_j(|\nabla u_{\epsilon_k}|)|\nabla u_{\epsilon_k}|^p)\,\mathrm{d}x\,\mathrm{d}t\nonumber\\
&\leqslant\lim_{k\to\infty}\iint_{Q_T}(1+|\nabla u_{\epsilon_k}|^p)\,\mathrm{d}x\,\mathrm{d}t\leqslant\widetilde C \label{align:ch08-8.57}
\end{align}
其中 $\widetilde C>0$ 与 $j$ 无关（回忆 $\nu_{(x,t)}(\mathbb{R}^N)=1$ 对几乎所有 $(x,t)\in Q_T$）。由于 $\lim_{j\to\infty}f_j(|y|)|y|^p=|y|^p$ 对所有 $y\in\mathbb{R}^N$，由\eqref{align:ch08-8.57} 与 Fatou 引理得
\begin{align*}
\iint_{Q_T}\mathrm{d}x\,\mathrm{d}t\int_{\mathbb{R}^N}(1+|y|^p)\,\mathrm{d}\nu_{(x,t)}(y)\leqslant\widetilde C,
\end{align*}
由\eqref{align:ch08-8.55} 结论即得。
\end{remark}

为后文引用，证明如下结果（见文献~\cite{ref100} 命题 6）。

\begin{proposition}[测度分量的分离等式]\label{proposition:ch08-分离等式}
设 $N=1$，$\Omega\equiv(a,b)$，\eqref{align:ch08-8.1a}--\eqref{align:ch08-8.1d} 成立。设 $(u,\nu)$ 为\cref{theorem:ch08-存在性} 给出的 (P) 的 Young 测度值解。设 $G_\phi$ 如\eqref{align:ch08-8.14}（$g\in C^1_c(\mathbb{R})$），$f\in C^1(\mathbb{R})\cap W^{1,\infty}(\mathbb{R})$。则对几乎所有 $(x,t)\in Q_T$，
\begin{align}
\int_{\mathbb{R}}f(\phi(y))G_\phi(y)\,\mathrm{d}\nu_{(x,t)}(y)
=\int_{\mathbb{R}}f(\phi(y))\,\mathrm{d}\nu_{(x,t)}(y)\int_{\mathbb{R}}G_\phi(y)\,\mathrm{d}\nu_{(x,t)}(y). \label{align:ch08-8.58}
\end{align}
\end{proposition}

为证\cref{proposition:ch08-分离等式}，需要一个技术性引理（证明见文献~\cite{ref83} 引理 2.1）。设 $\{\eta_i\}\subseteq H^1(\Omega)$ 为算子 $-\Delta$ 在 $\partial\Omega$ 上齐次 Neumann 条件下的特征函数序列，$\{\mu_i\}$ 为相应特征值序列。对任意 $\epsilon>0$，令 $\mathbb{N}_\epsilon\triangleq \{i\in\mathbb{N}\mid\mu_i\leqslant1/\epsilon\}$，并令 $P_\epsilon,Q_\epsilon:L^2(\Omega)\to H^1(\Omega)$、$P_\epsilon+Q_\epsilon=I$，为由下式定义的投影算子：
\begin{align}
P_\epsilon f\triangleq \sum_{i\in\mathbb{N}_\epsilon}f_i\eta_i,\quad
Q_\epsilon f\triangleq \sum_{i\in\mathbb{N}_\epsilon^c}f_i\eta_i,\quad
f_i\triangleq \int_\Omega f\eta_i\,\mathrm{d}x\ (f\in L^2(\Omega)). \label{align:ch08-8.59}
\end{align}

\begin{lemma}[投影估计]\label{lemma:ch08-投影估计}
设 $N=1$ 且\eqref{align:ch08-8.1a}--\eqref{align:ch08-8.1d} 成立。设 $u_\epsilon$ 为\cref{theorem:ch08-正则化适定} 给出的问题 $(P\epsilon)$ 的解，$v_\epsilon\triangleq \partial_x u_\epsilon$。则存在 $C>0$，使对所有 $\epsilon>0$，
\begin{align}
\|P_\epsilon\phi(v_\epsilon)\|_{L^2(0,T;H^1_0(\Omega))}
+\frac{1}{\sqrt\epsilon}\|Q_\epsilon\phi(v_\epsilon)\|_{L^2(Q_T)}\leqslant C. \label{align:ch08-8.60}
\end{align}
\end{lemma}

\begin{proof}[\cref{proposition:ch08-分离等式} 的证明]
设 $\{v_{\epsilon_k}\}$ 为\eqref{align:ch08-8.46} 中收敛的序列。在当前假设下，$f\circ\phi\in C_b(\mathbb{R})$，而 $|G_\phi(y)|\leqslant\|g\|_{L^\infty(\mathbb{R})}|y|$ 对所有 $y\in\mathbb{R}$。由\eqref{align:ch08-8.25} 序列 $\{v_{\epsilon_k}\}$ 在 $L^\infty(Q_T)$ 中有界因而在 $Q_T$ 中一致可积。于是由\cref{proposition:ch05-弱收敛} 在 $L^1(Q_T)$ 中有
\begin{align*}
f(\phi(v_{\epsilon_k}))\rightharpoonup(f\circ\phi)^*,\quad
G_\phi(v_{\epsilon_k})\rightharpoonup(G_\phi)^*,\quad
f(\phi(v_{\epsilon_k}))G_\phi(v_{\epsilon_k})\rightharpoonup((f\circ\phi)G_\phi)^*,
\end{align*}
其中
\begin{align*}
(f\circ\phi)^*\triangleq \int_{\mathbb{R}}f(\phi(y))\,\mathrm{d}\nu(y),\quad
(G_\phi)^*\triangleq \int_{\mathbb{R}}G_\phi(y)\,\mathrm{d}\nu(y),\quad
((f\circ\phi)G_\phi)^*\triangleq \int_{\mathbb{R}}f(\phi(y))G_\phi(y)\,\mathrm{d}\nu(y).
\end{align*}
于是若证明了
\begin{align}
f(\phi(v_{\epsilon_k}))G_\phi(v_{\epsilon_k})\rightharpoonup(f\circ\phi)^*(G_\phi)^*\ \text{在 }L^1(Q_T), \label{align:ch08-8.61}
\end{align}
结论即得。令 $F\triangleq f\circ\phi$，$F_{\epsilon_k}\triangleq f(P_{\epsilon_k}\phi(v_{\epsilon_k}))$，其中 $P_{\epsilon_k}$ 为\eqref{align:ch08-8.59} 中取 $\epsilon=\epsilon_k$ 的投影算子。由假设 $\|f'\|_{L^\infty(\mathbb{R})}<\infty$，有
\begin{align}
\|F_{\epsilon_k}-F(v_{\epsilon_k})\|_{L^2(Q_T)}
=\|f(P_{\epsilon_k}\phi(v_{\epsilon_k}))-f(\phi(v_{\epsilon_k}))\|_{L^2(Q_T)}
\leqslant\|f'\|_{L^\infty(\mathbb{R})}\|Q_{\epsilon_k}\phi(v_{\epsilon_k})\|_{L^2(Q_T)}. \label{align:ch08-8.62}
\end{align}
另一方面，由\eqref{align:ch08-8.25} 存在 $C>0$ 使 $\|G_\phi(v_{\epsilon_k})\|_{L^2(Q_T)}\leqslant C$。由\eqref{align:ch08-8.60} 与\eqref{align:ch08-8.62} 得
\begin{align*}
\lim_{k\to\infty}\|[F_{\epsilon_k}-F(v_{\epsilon_k})]G_\phi(v_{\epsilon_k})\|_{L^2(Q_T)}=0.
\end{align*}
因此，为证\eqref{align:ch08-8.61}，只需对每个 $\zeta\in C^1_c(Q_T)$ 证明
\begin{align}
\lim_{k\to\infty}\iint_{Q_T}F_{\epsilon_k}G_\phi(v_{\epsilon_k})\zeta\,\mathrm{d}x\,\mathrm{d}t
=\iint_{Q_T}(f\circ\phi)^*(G_\phi)^*\zeta\,\mathrm{d}x\,\mathrm{d}t. \label{align:ch08-8.63}
\end{align}
令 $\Gamma_{\epsilon_k}(x,t)\triangleq \int_a^x G_\phi(v_{\epsilon_k}(z,t))\,\mathrm{d}z$（$(x,t)\in Q_T$）。则
\begin{align}
\iint_{Q_T}F_{\epsilon_k}G_\phi(v_{\epsilon_k})\zeta\,\mathrm{d}x\,\mathrm{d}t
=-\iint_{Q_T}\partial_x(F_{\epsilon_k}\zeta)\Gamma_{\epsilon_k}\,\mathrm{d}x\,\mathrm{d}t. \label{align:ch08-8.64}
\end{align}
由于 $\|f'\|_{L^\infty(\mathbb{R})}<\infty$，由不等式\eqref{align:ch08-8.60} 存在 $\bar C>0$ 使
\begin{align*}
\|\partial_x F_{\epsilon_k}\|_{L^2(Q_T)}
\leqslant\|f'\|_{L^\infty(\mathbb{R})}\|\partial_x[P_{\epsilon_k}[\phi(v_{\epsilon_k})]]\|_{L^2(Q_T)}
\leqslant\bar C\|f'\|_{L^\infty(\mathbb{R})}.
\end{align*}
由于 $f\in L^\infty(\mathbb{R})$，序列 $\{F_{\epsilon_k}\}$ 在 $L^\infty(Q_T)$ 因而在 $L^2(Q_T)$ 中有界；由上式它在 $L^2(0,T;H^1(\Omega))$ 中亦有界。于是存在 $\bar F\in L^2(0,T;H^1(\Omega))$ 使 $F_{\epsilon_k}\rightharpoonup\bar F$ 在 $L^2(0,T;H^1(\Omega))$。由\eqref{align:ch08-8.62} 这蕴含 $F(v_{\epsilon_k})\rightharpoonup\bar F$ 在 $L^2(Q_T)$，故 $(f\circ\phi)^*=\bar F$ 在 $Q_T$。因而 $(f\circ\phi)^*\in L^2(0,T;H^1(\Omega))$，且
\begin{align}
\partial_x(F_{\epsilon_k}\zeta)\rightharpoonup\partial_x[(f\circ\phi)^*\zeta]\ \text{在 }L^2(Q_T). \label{align:ch08-8.65}
\end{align}
另一方面，由于序列 $\{G_\phi(v_{\epsilon_k})\}$ 在 $L^\infty(Q_T)$ 中有界，存在 $G\in L^\infty(Q_T)$ 使（可能抽子列，不重新标号）$G_\phi(v_{\epsilon_k})\stackrel{*}{\rightharpoonup}G$ 在 $L^\infty(Q_T)$。由于 $G_\phi(v_{\epsilon_k})\rightharpoonup(G_\phi)^*$ 在 $L^1(Q_T)$，有 $G=(G_\phi)^*$。于是例行证明给出对几乎所有 $(x,t)\in Q_T$，
\begin{align}
\Gamma_{\epsilon_k}(x,t)\to\Gamma^*(x,t)\triangleq \int_a^x(G_\phi)^*(z,t)\,\mathrm{d}z \label{align:ch08-8.66}
\end{align}
（细节见文献~\cite{ref100} 命题 5 与 6）。由于族 $\{\Gamma_{\epsilon_k}\}$ 在 $L^\infty(Q_T)$ 中有界，由\eqref{align:ch08-8.66} 与控制收敛定理得
\begin{align}
\Gamma_{\epsilon_k}\to\Gamma^*\ \text{在 }L^2(Q_T). \label{align:ch08-8.67}
\end{align}
由\eqref{align:ch08-8.65} 与\eqref{align:ch08-8.67}，当 $k\to\infty$ 时等式\eqref{align:ch08-8.64} 的右端收敛于
\begin{align*}
-\iint_{Q_T}\partial_x[(f\circ\phi)^*\zeta]\Gamma^*\,\mathrm{d}x\,\mathrm{d}t
=\iint_{Q_T}(f\circ\phi)^*(G_\phi)^*\zeta\,\mathrm{d}x\,\mathrm{d}t.
\end{align*}
这就证明了\eqref{align:ch08-8.63}，结论即得。
\end{proof}




\end{document}